\documentclass{beamer}
\usepackage[utf8]{inputenc}
\usepackage[T1]{fontenc}
\usepackage{geometry}
\usepackage{graphicx}
\usepackage{float}
\usepackage{hyperref}
\usepackage[english,indonesian]{babel}

\title{Optimasi Pemrosesan Template LaTeX Menggunakan Pendekatan Berbasis AI}
\author{Academic Support Agent}
\date{\today}

\begin{document}
\maketitle

\begin{abstract}

\end{abstract}

\chapter{PENDAHULUAN}


\chapter{TINJAUAN PUSTAKA}


% Original content (commented out):
% % XeLaTeX document
\documentclass[aspectratio=169, 11pt]{beamer} % для экранов 16:9 оставить

% \documentclass[11pt]{beamer} % для экранов 4:3 оставить

% Редактируем: конфигурация, личные настройки: имя, название предмета и пр. для титульной страницы и метаданных документа здесь
\include{config}

% Не редактируем: используемые пакеты
\include{settings/packages}

% Не редактируем: параметры используемых пакетов и не только
\include{settings/preferences}

% водяной знак для обозначения статуса документа
% \newwatermark[allpages,color=red!5,angle=45,scale=3,xpos=0,ypos=0]{DRAFT}

% настройка метаданных документа
\institute[\university]{\university \\ \faculty}
\title[\docname]{\docname}
\author[\studentname]{\small Выполнил: ст. гр. \group \ \studentname\\[1mm]{\small Руководитель: \tutorname}}
\date{\the\year} 

\include{templates/titlepage}

\begin{document}

% Редактируем: всё остальное: вступление, др. этапы, заключение, приложение
\include{content}

% Не редактируем: Страница библиографии (формируется автоматически из книжек, указанных в refs.bib и пометок \cite{имя_источника} в тексте)
\include{templates/bib-page}

\end{document}


\end{document}